\section{Finiteness and Growth}

The mesh is finite at every beat but always expanding. There is never an actual infinity anywhere in the model, only a finite region that grows without end. This section sets out what that means, why no completed infinity is ever realized, and how small a finite chunk of space has to be before it already looks like today's space. The answer to the last question turns out to be surprisingly small, and it was measured.

\begin{principle}[The receding edge]
The ever-receding edge of the mesh is the present. Arriving would be the end of time.
\end{principle}

\subsection{Finite at any beat, ever-growing}

The base space is the $\{3,4,3,4\}$ hyperbolic honeycomb. Its extent is the property that matters here. The mesh is neither an infinite static space laid out once and for all, nor a closed torus that wraps back on itself. It is something stranger and more dynamic. At any single beat it is a finite region with a real boundary, and at the next beat it is a larger finite region with a boundary further out.

\begin{table}[ht]
\centering
\begin{tabular}{@{}ll@{}}
\toprule
property & the mesh \\
\midrule
infinite static space & no \\
closed torus & no \\
finite at any beat, ever-growing & yes \\
\bottomrule
\end{tabular}
\caption{The extent of the mesh. It is finite at every beat and grows without end.}
\label{tab:mesh-extent}
\end{table}

The real boundary of the mesh is the \emph{wake}. The wake is the growing edge, the frontier where the lattice ends. Each beat the wake moves outward as the lattice grows, and new empty docks are born at peace, with every site holding the tone $0$, at that frontier. These fresh docks are the open future arriving.

\begin{definition}[The wake]
The \emph{wake} is the outer boundary of the mesh at a given beat. It is the ever-expanding frontier where new docks wake into existence, each born at peace. The wake moves outward by one shell each beat, and the docks behind it are never revisited from beyond.
\end{definition}

The arrow of time and the radiation bath both depend on this growth. The wake is a low-tone source feeding fresh, peaceful docks into the mesh, and what falls behind the wake is carried away and never returns. The growth is therefore not a cosmetic add-on. It is the engine of the arrow.

\begin{principle}[No actual infinity in the mesh]
The mesh never contains an actual infinity. It is finite at every beat, and it never stops growing.
\end{principle}

\subsection{The complete cusp is never reached, and need not be}

Physical space is the cubic cusp of $\{3,4,3,4\}$, the flat $\{4,3,4\}$ region that opens out toward a cusp at infinity. The \emph{complete} cusp is the infinite limit, the perfect flat $D_4$ space that the patch is forever flowing toward.

That limit is never realized in a finite mesh. At any finite beat the mesh has only a finite chunk of $\{4,3,4\}$, not the whole of it. Physical space is therefore never the complete cusp. It is an ever-growing finite piece of it, a horosphere patch at some finite cusp-height, slightly curved, with a real wake for a boundary.

This is exactly the situation of a finite crystal that uses the infinite-lattice idealization. A real crystal is a finite lump of atoms, yet its physics is computed with the infinite periodic lattice, because local physics only ever sees the local lattice. The boundary is far away and irrelevant to what happens in the middle. The mesh works the same way.

\begin{table}[ht]
\centering
\begin{tabular}{@{}ll@{}}
\toprule
at a finite beat & in the $t \to \infty$ limit \\
\midrule
a finite horosphere patch & the infinite perfect flat $D_4$ space \\
at finite cusp-height, slightly curved & perfectly flat \\
with a real wake & no boundary \\
\bottomrule
\end{tabular}
\caption{The finite chunk of physical space at any beat, against its infinite attractor.}
\label{tab:patch-vs-limit}
\end{table}

Eternal expansion is the patch flowing ever deeper toward the cusp. It grows and flattens a little each beat, forever approaching the ideal infinite flat $D_4$ space but never reaching it. So both statements hold at once, each at its proper place. Finite-and-growing is the reality at every beat. The infinite perfect flat cusp is the attractor that the reality forever approaches.

\begin{result}[Never-reaching is the open future]
The never-reaching of the complete cusp is the eternal expansion and the open future. The perfectly flat infinite cusp is a limit, not a state the lattice is ever in at a finite beat. Arriving would be the end of time.
\end{result}

\subsection{Convergence is fast}

How big must the finite, growing chunk of cusp be before it looks like today's space? The honest way to answer is to build it and measure. The cusp is the cheap $\{4,3,4\}$ cubic lattice, so finite cubic chunks of growing side $L$ are inexpensive to construct and probe. For each side $L$ two quantities were measured, the spectral dimension seen by a random walk and the exponent of the emergent gravitational law.

\begin{table}[ht]
\centering
\begin{tabular}{@{}rll@{}}
\toprule
side $L$ & spectral dimension & gravity exponent \\
\midrule
5  & 2.58 & 2.06 \\
7  & 2.88 & 1.61 \\
11 & 2.99 & 1.34 \\
17 & 2.99 & 1.40 (converged) \\
25 & 2.99 & 1.42 \\
35 & 2.99 & 1.43 \\
\bottomrule
\end{tabular}
\caption{Measured spectral dimension and gravity exponent of cubic cusp chunks of side $L$.}
\label{tab:convergence}
\end{table}

Two readings come straight off the table. The spectral dimension reaches $3$ by side $L \approx 11$ cells, so the space already walks like three-dimensional space at that scale. Local physics becomes continuum-like, meaning $L$-independent, by side $L \approx 17$ cells. A cube of side $17$ holds about $5000$ cells.

\begin{result}[Continuum physics arrives in a few beats]
Clean, $L$-independent local physics arrives by side $L \approx 17$, about $5000$ cells. Mapping size to growth through $a(t) \sim e^{H t}$, that is only a few beats of expansion. This is a measured emergent result on the cusp geometry.
\end{result}

\subsection{The observed universe is vastly beyond sufficiency}

The observed universe is about $10^{61}$ cells across. Set that against the thresholds just measured, and the comparison is stark.

\begin{table}[ht]
\centering
\begin{tabular}{@{}lll@{}}
\toprule
scale & cells & regime \\
\midrule
spectral dimension $3$ & $\sim 11$ (side $L$) & reached almost at once \\
continuum local physics & $\sim 17$ (side $L$), $\sim 5000$ cells & a few beats \\
observed universe & $\sim 10^{61}$ & astronomically deep in the continuum regime \\
\bottomrule
\end{tabular}
\caption{The observed universe against the convergence thresholds.}
\label{tab:universe-scale}
\end{table}

So finite-size and edge effects at any observable scale are utterly negligible. The wake is unimaginably far from anything we could probe.

A note on the cheap computation is in order, since these numbers might seem to require building astronomically many cells. They do not. The cusp's intrinsic structure is the $\mathbb{Z}^3$ cubic lattice, of order $L^3$ cells, so the boxes are built directly at any size, with no expensive hyperbolic horoball to grow first. The $10^{61}$ figure is reached by analytic extrapolation, $\mathbb{Z}^3$ tending toward isotropic three-dimensional elasticity as a power of $1/L$, not by constructing that many cells.

\subsection{The refining universe}

A picture of cosmic history falls out of all this. The universe expanded fast to a rough geometry, a small chunk that was already roughly three-dimensional but not yet clean. Then it refined toward perfect three-dimensional space as it grew. Flatness and isotropy improved with every beat of expansion.

\begin{table}[ht]
\centering
\begin{tabular}{@{}ll@{}}
\toprule
stage & geometry \\
\midrule
early, small chunk & rough, sub-three-dimensional, anisotropic \\
growing & refining toward clean three dimensions, flatness and isotropy improving \\
late, vast chunk & nearly perfect three dimensions, approaching but never reaching the infinite-cusp limit \\
\bottomrule
\end{tabular}
\caption{Cosmic history as a refining chunk of cusp.}
\label{tab:refining}
\end{table}

\begin{principle}[The steady burn]
The fire reached a steady, clean burn almost at once, and has been steadying ever since.
\end{principle}

This matches observed cosmology. Space is measured flat to about $0.4$ percent and the universe has a finite age, both consistent with a finite chunk that grew fast and has been refining toward flatness. One thing does not vanish with growth. The fixed cubic anisotropy from the finite octahedral cusp symmetry is a permanent feature, not a finite-size artifact. It is a tiny imperfection that no amount of further growth removes, since it is set by the symmetry of the cusp rather than by the size of the patch.

\subsection{Whether it began is open, that it grows is committed}

The theory takes a clear stance on the future and a deliberately open stance on the past. It commits to endless growth. It declines to commit to a first beat.

\begin{table}[ht]
\centering
\begin{tabular}{@{}ll@{}}
\toprule
question & stance \\
\midrule
does the mesh grow without end & committed, yes \\
was there a first beat & left open \\
\bottomrule
\end{tabular}
\caption{The committed future and the open past.}
\label{tab:open-past}
\end{table}

The forward growth is load-bearing. It is the source of the irreversibility that every self requires, the low-tone frontier against which the arrow of time is defined. The wake carries the past away, which is precisely why the present can heal and selves can persist. Without ongoing growth there is no arrow and no bath.

\begin{result}[Growth committed, origin open]
That the mesh grows without end is committed. Whether it began is left open. The dynamics, the arrow, the bath, and the self all live in the ongoing growth, not in any origin event. The origin is a separate question on which the model stays agnostic.
\end{result}

The theory does not need a first beat to function. Everything load-bearing happens in the growing, not in any starting. Whether the growth had a beginning is a further question that can be left to one side without affecting any of the physics.

\subsection{Never an actual infinity}

The finiteness stance is easy to misread, so it pays to state it plainly and in one place. The mesh is finite at every beat, a real bounded region with a wake. The mesh grows without end, exponentially, each shell of new docks bigger than everything inside it. The infinite flat cusp is a limit, an attractor, never an actual state of the lattice. There is never a completed infinity anywhere in the model.

\begin{principle}[Potential, not actual]
This is a potential infinity, not an actual one. The mesh is always growing toward the ideal, never arriving. Finite at every beat. Growing without end. Never an actual infinity.
\end{principle}

\subsection{Why it matters}

A finite chunk of cusp that sits far from any infinity is more than sufficient to look like today's space. The continuum regime arrives in a few beats, around $5000$ cells, and the real universe is some $10^{61}$ cells across, astronomically beyond that threshold. The model needs no actual infinity to reproduce flat three-dimensional physics. It needs only a finite, ever-growing chunk, which it has at every beat.

The growth is not a separate ingredient bolted on. It is the same growth that is time. It is the same growth that feeds the bath. It is the same growth that gives the arrow. Finiteness and growth are one fact seen from three sides, and that single fact is enough to deliver flat space, a forward arrow, and an open future together.
